
\section{Concerning Lay Activity}

Editorial article in “Foldebladet” for April 3, 1907.

The great Joint Committee, which is negotiating concerning a union of the church bodies, also dealt with “lay activity” at its meeting last autumn. The result of the work is set forth in 7 propositions, which, it is said, were unanimously adopted by the representatives from the different bodies. To each of the 7 propositions there is added one or more passages of Scripture, or else statements from the Lutheran confessional writings. The connection, however, between the passages of Scripture and the propositions is so unclear that explanation is needed in order to make them fit together.

Neither do the propositions speak of “lay activity” alone, nor even first and foremost. They speak of the congregation, the Means of Grace, the preaching office, the gifts of grace, the universal priesthood, and finally perhaps of “lay activity.” But although all these things are spoken of, there is nevertheless no explanation either of what is meant by the congregation or by “lay activity.” Indeed, we say that perhaps they speak of “lay activity”; for it is not mentioned by name in any of the propositions. Properly speaking, we infer that proposition 5 speaks of “lay activity,” because there are certain expressions which somewhat call this matter to mind, and because such great care is used to get it properly controlled.

Proposition 5 reads as follows:

“God also wills that the particular gifts of grace which He may have given to certain persons in the congregation, and which are therefore the property of the congregation, shall be put to use by the congregation, thus also in particular the gift of grace of being able to proclaim the Word of God to an assembly. According to the Word of God, the one who has received this gift of grace shall be called to this service in the congregation by its invitation or consent, and shall use the gift under the supervision of the congregation or of those whom the congregation has called to exercise supervision over the work in the congregation, as the Word of God and the Confession of our Church say.”

Nothing is said as to what is meant by the congregation, either in this or in any of the other propositions; only we read in proposition 1 that “God has given the congregation, also the individual local congregation, the Means of Grace,” etc. Presumably, then, by the congregation is meant the purely outward organization which in each place bears the name of congregation, whether it deserves the name or not, whether it is dead or living, awake or asleep.

If we look at the proposition with this assumption, then we find the following assertions, some with and some without testimony that they are “according to the Word of God”:

* 1. When certain persons in the congregation have received a gift of grace, then for that reason this is “the property of the congregation.”
* 2. When someone has received “the gift of proclaiming the Word of God to an assembly,” he shall be called to this service in the congregation by its invitation or consent.
* 3. When he has thus been called, he shall use the gift under the supervision of the congregation or of those whom the congregation has called to exercise supervision over the work in the congregation.

The first of these assertions stands on its own worth; but the second is said to be “according to the Word of God”; the third, “as the Word of God and the Confession of our Church teach.”

We have supposed that in this proposition 5 the Joint Committee is thinking of “lay activity”; for the quotation marks that are used must here denote that the word is taken from common usage, which has made the word so well known among us; and the “lay activity” commonly known among us may perhaps be designated as a “gift of grace to proclaim the Word of God in an assembly,” without our thereby acknowledging that the explanation is exact or sufficient.

Now in proposition 5 it is stated that all gifts of grace which God has given to certain persons in the congregation are therefore the property of the congregation. By what follows, this means that the congregation can and shall determine whether these gifts are to be used, and how they are to be used, of course under responsibility to God, but with such power and authority that no one has the right to use the gift—at any rate not the one of speaking in an assembly—without the congregation’s permission.

Here we must ask: Where is it written that the individual Christian’s gift of grace is the property of the congregation? “All things are yours,” says Paul, but that is surely meant of Christians, not of an outward organization of people concerning whose Christianity only this can be said, that they are entered in a congregational register? Christians can say that all is ours because they are Christ’s and Christ is God’s; but is it thereby said that an outward organization of Christians and non-Christians can rule and command over all things in Heaven and on earth and in particular can determine whether the Christian shall use his gift of grace or not? Can a congregational decision remove the responsibility from those who have received the gift, so that he need not use it because the congregation does not invite him to do so? The Word of God says that the individual Christian has responsibility for his gift of grace or his talent. The Word of God says that only the one who has interest from his pound is allowed to enter into his Lord’s joy, and that the one who has no interest, because he has hidden his pound safely in the earth, is cast out into the outer darkness, where there is weeping and gnashing of teeth. See how altogether differently that language sounds from proposition 5 of the Committee, and it gives no excuse for the slothful flesh.

No, it is said, the Committee has meant nothing of the sort except concerning the gift of proclaiming the Word of God in an assembly. “Lay activity” in the Haugean sense, then, is the only gift which the congregation shall acknowledge by invitation or consent, and the one who has received it shall conform himself accordingly. If he receives an invitation, he shall speak; if he receives no invitation, he shall keep silence in an assembly.

If now this is the meaning, does it make the matter better? Or is it written anywhere that the one who has received this gift is free of his responsibility when the congregation does not invite him?

Suppose now, then, that the congregation sleeps, and God in His grace sends a crying voice to awaken it? Or suppose that the congregation is dead, even though it has the name of being alive; shall then the one who was sent even to awaken the spiritually dead keep silence until the dead invite the living one to speak?

No, neither the first assertion, that the gifts of grace are the property of the congregation, nor the second assertion, that no one may bear witness in an assembly until he is invited by the congregation, stands “according to the Word of God,” which much rather makes the duty of witness depend not on a congregational call but simply and straightforwardly on the personal Christian calling. Whether the testimony is borne to one person or “in an assembly” makes no difference according to Scripture. Let the one who has the gift to speak to the individual speak to the individual, and let the one who has the gift to speak in an assembly speak in an assembly.

What does one mean, in any case, by “an assembly”? Take note that it does not say the congregation’s assembly or the congregation; it says only “an assembly.” How many make up an assembly, where and when must people meet together in order that it may become an assembly? Anyone can understand that such indefinite words can be used by a false congregation or pastor to suppress the proclamation of the Word of God and all testimony from individual Christians in and outside the congregation. The French law in Madagascar determines that a Christian evangelist or missionary may outside the church building speak the Word of God to two but not to three; for three are an assembly; at home in France the law says that a Christian preacher may speak the Word of God outside the church building to 19 but not to 20 persons; for 20 persons are an assembly. Are not the enemies of living Christianity of the same sort in America as in France? And are there not many, many congregations where the majority is against lay activity?

We ought not to frame our rules so that the enemies of revival get so good a grip to hold on to as this. Precisely by such means Hauge in Norway was made a criminal, because without men’s call he spoke the Word of God in assemblies. Shall we return to such rules, which here in our land may indeed not bring a man into fines and prison, but which can nonetheless shut his mouth, so that he keeps silence where he ought to speak?

And the last assertion in proposition 5, that lay activity shall be exercised under the supervision of the congregation or of those whom the congregation has chosen for that purpose, plainly goes in the same direction. For this supervision must in the connection mean that congregation or pastor or congregational council has the right in turn again to forbid the one who has the gift to use it, if they so wish. For it comes to the same thing, namely this: that God can give gift and call, but men, because they have formed an outward congregational order, can forbid the use of the gift.

To all this the Committee can say that it has a proposition 6 which declares that in such a case “the congregation suffers harm.” But that does not cover the matter; for the question is this: Can the one who has received the gift refrain from using it for any reason whatsoever without himself suffering harm? And it is a harm which according to the Word of God consists in this, that “his hands and feet are bound, and he is cast out into the outer darkness, where there is weeping and gnashing of teeth.”

Georg Sverdrup
